\documentclass[11pt]{article}
\usepackage[a4paper,margin=2.6cm]{geometry}
\usepackage[T1]{fontenc}
\usepackage[utf8]{inputenc}
\usepackage{lmodern,microtype,amsmath,amssymb,booktabs,array,enumitem,graphicx}
\usepackage[backend=biber,style=authoryear,natbib=true,maxcitenames=2,maxbibnames=99]{biblatex}
\addbibresource{references.bib}
\usepackage[hidelinks]{hyperref}
\usepackage[nameinlink,noabbrev]{cleveref}
\newcommand{\Prob}{\mathbb{P}}
\newcommand{\KL}{D_{\mathrm{KL}}}

\title{When Confidence Is Not Information:\\Failure Modes of Abstention in Ambient Behavioural Sensing}
\author{Diogo Ribeiro\\\small\textit{Faculty of Media Arts and Design, Technical University of Porto}}
\date{}

\begin{document}
\maketitle

\begin{abstract}
Probabilistic behavioural-sensing systems often use posterior confidence as a proxy for whether a prediction is trustworthy. We test that proxy in an interpretable smart-home state-inference system using posterior confidence, posterior margin, normalised entropy, interval-level evidence strength, and per-update information gain. Correct and incorrect predictions are compared within each household using a common-language AUC and households are weighted equally. Across 22 CASAS development homes, median correctness AUC is 0.519 for confidence, 0.518 for posterior margin, and 0.514 for normalised entropy. Evidence strength and information gain are inverted, with median AUCs of 0.372 and 0.388. Thus all three directional hypotheses fixed before execution fail: information gain does not exceed 0.5, does not outperform confidence, and evidence strength does not carry positive correctness information. The result rules out a simple replacement of posterior confidence by update-level information gain as an abstention statistic on this development panel. No new threshold or deployment rule is selected.
\end{abstract}

\section{Introduction}
An abstaining classifier claims not only to rank hypotheses but also to recognise when evidence is insufficient for a decision. In ambient behavioural sensing this is difficult because quiet sensors, missing data, sparse observations, and persistent latent-state dynamics can all produce concentrated posteriors without reliable predictions.

The system studied here already abstains using posterior confidence and sensor completeness. Earlier development analyses of the same system \citep{ribeiro2026toolkit} showed that this is a weak basis for abstention. Over 60,948 scored steps from five CASAS homes, mean confidence was 0.832 on correct and 0.758 on incorrect predictions, and the 0.95--1.00 confidence band, which held 39\% of those steps, was less accurate (0.561) than the 0.85--0.95 band (0.653). In confirmatory simulation, mean abstention rose only from $9.7\times10^{-6}$ to $5.3\times10^{-5}$ as random missingness rose from 0\% to 40\%. This paper asks whether another internal diagnostic ranks prediction correctness more usefully.

\section{Related work}
Classification with a reject option goes back to \citet{chow1970reject}, who characterised the optimal trade-off between error and rejection when posterior probabilities are known. Selective prediction studies the same trade-off between coverage and error without assuming that a fixed posterior threshold is meaningful \citep{geifman2017selective,geifman2019selectivenet}. Calibration is related but distinct: a model may be calibrated while ranking failures poorly, or rank failures while remaining miscalibrated \citep{guo2017calibration}. Misclassification detection evaluates a confidence score by how well it ranks correct above incorrect predictions, usually summarised by the area under the ROC curve \citep{hendrycks2017baseline}; the within-home AUC used here is that estimand computed per household. Explicit failure-prediction work likewise shows that maximum class probability can be a weak failure signal \citep{corbiere2019failure}. Our analysis remains model-internal and non-learned: we compare quantities already present in the probabilistic trace before introducing any new reject policy.

\section{Diagnostics}
Let $p_{t|t-1}$ be the predicted latent-state distribution before the observations at time $t$ and $p_t$ the posterior after their likelihoods are applied.

Posterior confidence is
\[
C_t=\max_z p_t(z),
\]
and posterior margin is the gap between the two largest posterior probabilities. Normalised entropy is
\[
\widetilde H_t=-\frac{1}{\log|\mathcal S|}\sum_z p_t(z)\log p_t(z).
\]
If $s_{jt}$ is the log-likelihood support contributed by informative sensor $j$, interval-level evidence strength is
\[
E_t=\frac{1}{|J_t|}\sum_{j\in J_t}|s_{jt}|.
\]
Per-update information gain is
\[
I_t=\KL\!\left(p_t\,\|\,p_{t|t-1}\right).
\]
Unlike confidence, margin, and entropy, $I_t$ measures how much the current observation interval moved the predicted belief.

\section{Empirical design}
The analysis uses the 22 CASAS \texttt{hh} homes previously used for development and failure analysis, taken from the published labelled CASAS archive \citep{cook2013casas,cook2025casasdata}. By the resident counts in that archive's metadata, 20 are single-resident and two, hh107 and hh121, are two-resident recordings; both are retained because panel membership was fixed before this analysis. These homes are not independent confirmatory data. The frozen 43-home primary external cohort from the 0.3.0 validation is not rescored here.

Each recording is run through the default inference path, with the circadian profile disabled, on a five-minute grid with timestamps read in the America/Los\_Angeles time zone. Only positions with a mapped behavioural label are scored, 242,681 positions in total, and a prediction is correct when the most probable state equals the label.

For a diagnostic $X$, let $X^+$ and $X^-$ be values on correct and incorrect predictions. For confidence, margin, evidence strength, and information gain,
\[
A_X=\Prob(X^+>X^-)+\frac12\Prob(X^+=X^-).
\]
For entropy the favourable direction is reversed. Thus $A=0.5$ means no separation, $A>0.5$ favours correct predictions, and $A<0.5$ indicates inversion.

The household is the aggregation unit. We compute $A_{X,h}$ within each household and report the median across households, so long recordings do not dominate short ones.

\subsection{Pre-result contract}
\label{sec:contract}
The hypotheses and their interpretation were committed to the project repository before the analysis was first executed, in commit \texttt{34f9797} on 13 September 2026. \Cref{app:contract} reproduces that text verbatim. It fixed three directional hypotheses:
\begin{enumerate}[label=H\arabic*.,leftmargin=*]
\item $\operatorname{median}_h(A_{I,h})>0.5$;
\item $\operatorname{median}_h(A_{I,h})>\operatorname{median}_h(A_{C,h})$;
\item $\operatorname{median}_h(A_{E,h})>0.5$.
\end{enumerate}
It also fixed that no significance threshold, equivalence margin, or post-hoc cutoff would be introduced after the panel was run, and that no abstention threshold, coverage target, or deployment rule would be selected from these 22 homes.

The panel runner had been committed 84 minutes before the contract. Between the contract commit and the executed run, the analysis code changed only to verify each source recording against its stored digest, and the package version number changed. The earliest recorded execution is a continuous-integration run on 14 September 2026. Local executions are not logged, so the repository history alone cannot rule out an earlier unrecorded run.

\section{Results}
\begin{table}[ht]
\centering
\caption{Within-home correctness AUC on the 22-home CASAS development panel. Homes counts the households in which the diagnostic is defined; the last column counts those with $A>0.5$. No household has $A=0.5$ exactly.}
\label{tab:separation}
\begin{tabular}{lcccc}
\toprule
Diagnostic & Median AUC & Homes & Range & Homes with $A>0.5$ \\
\midrule
Confidence & 0.519 & 22 & 0.164--0.707 & 12 \\
Posterior margin & 0.518 & 22 & 0.166--0.693 & 12 \\
Normalised entropy & 0.514 & 22 & 0.160--0.714 & 12 \\
Evidence strength & 0.372 & 22 & 0.103--0.740 & 2 \\
Information gain & 0.388 & 22 & 0.111--0.911 & 4 \\
\bottomrule
\end{tabular}
\end{table}

\begin{figure}[ht]
\centering
\includegraphics[width=0.92\linewidth]{uncertainty_panel_auc.pdf}
\caption{Within-home correctness AUC for all 22 development homes and all five diagnostics. Each point is one household. The dashed horizontal line marks the fixed no-separation reference $A=0.5$. No fitted trend, threshold line, post-hoc filtering, or household weighting is applied.}
\label{fig:uncertainty-panel}
\end{figure}

Median within-home accuracy is 0.489, ranging from 0.039 to 0.586.

Confidence, margin, and entropy are close to the no-separation value at the median, with confidence the largest of the three at 0.519. They are not uniformly uninformative: each exceeds 0.5 in 12 of 22 homes, and confidence ranges from 0.164 to 0.707 across homes.

The update-level diagnostics are consistently inverted. Information gain is below 0.5 in 18 of 22 homes, with median 0.388, and evidence strength in 20 of 22, with median 0.372: larger updates and larger evidence magnitude are associated with incorrect predictions more often than with correct ones in the common-language ranking sense. Their medians lie 0.112 and 0.128 below 0.5, against 0.019 above it for confidence, so at the median they rank correctness more strongly than confidence does, but with the opposite sign to the one hypothesised.

Consequently, H1 fails because $0.388<0.5$, H2 fails because $0.388<0.519$, and H3 fails because $0.372<0.5$. All five diagnostics are defined in all 22 homes.

One home, hh124, is anomalous. Only 666 of its 17,064 scored positions are correct (3.9\%), its source recording is 0.53\,MB against 2.0--11.9\,MB for the other 21 homes, and it holds the lowest AUC for confidence, margin, and entropy and the highest for evidence strength and information gain. Its cause has not been investigated. As a post-hoc check, excluding hh124, the two two-resident homes, or all three moves no median by more than 0.021 and leaves H1, H2, and H3 failing.

\section{Interpretation}
The executed result follows the negative branch of the pre-result contract.

First, information gain is not a replacement for confidence. The original motivation was that a posterior can remain concentrated while the current interval adds little new information, making $I_t$ a plausible alternative. The real panel rejects that proposal: larger information gain is associated with incorrect predictions more often than with correct ones at the household median.

Second, evidence strength is also inverted. The problem is therefore not simply that useful raw evidence is transformed badly by the Bayesian update; the magnitude of the applied evidence itself does not provide a favourable correctness ranking under the current formulation.

Third, confidence, margin, and entropy are only weakly informative. None of the five existing scalar diagnostics emerges as a credible abstention statistic. The failure is therefore deeper than a badly chosen confidence threshold: the present observation and inference formulation can produce strong updates in the wrong direction while posterior concentration itself carries little correctness information.

This does not imply that the sensors contain no useful information. In earlier development analyses, a gradient-boosted classifier given per-room event counts for the current and three preceding steps together with time of day, trained on 11 of these homes and scored on the other 11, reached balanced accuracy 0.607, against a median of 0.420 for the inference pipeline across the 22 homes \citep{ribeiro2026toolkit}. The present claim is narrower: these five internal uncertainty summaries do not reliably rank prediction correctness across the development homes.

No post-hoc abstention threshold is selected. Any future selective-prediction policy first requires a candidate score with defensible ranking behaviour and then an independent validation design.

\section{Limitations}
The 22 homes are development-stage evidence and were repeatedly inspected during earlier failure analysis. The results should not be read as confirmatory external performance. All homes come from the same CASAS research ecosystem.

Information gain measures how much the model updates itself; it is not an external measure of information in the Shannon-channel sense. The common-language AUC is a ranking measure, not a calibration measure. The inversion of evidence strength and information gain is therefore a property of these diagnostics under this observation model, inference formulation, state ontology, and development panel, not a universal property of those quantities.

\section{Reproducibility}
The analysis is implemented by \texttt{scripts/analyse\_uncertainty\_panel.py}. It reads the frozen 22-home development-panel membership, verifies each source recording against its stored byte size and SHA-256 digest, evaluates each home using the default inference path, and records per-home diagnostics and the equal-household summary.

The executed workflow artifact is frozen at \texttt{artifacts/paper1/uncertainty\_panel.json}. It is the source of truth for every value in \cref{tab:separation} and for the accuracy, outlier, and exclusion values in the text, which are computed from its per-home entries. The pre-specified household-distribution figure in \cref{fig:uncertainty-panel} is rendered deterministically from the same artifact by \texttt{scripts/plot\_uncertainty\_panel.py}.

The workflow does not pin dependency versions, so results are reproducible to numerical tolerance rather than bit for bit. An earlier execution of the same analysis code on the same day differs from the frozen artifact by at most $2\times10^{-6}$ in any per-home value and yields identical medians.

\section{Conclusion}
Posterior confidence is not a strong correctness-ranking signal in this ambient behavioural-sensing system, but the natural update-level alternatives perform worse rather than better. Across 22 development homes, confidence, margin, and normalised entropy have median correctness AUCs only slightly above 0.5, while evidence strength and information gain are inverted at 0.3719 and 0.3875. All three directional hypotheses specified before execution fail.

The practical implication is narrow but important: the abstention problem cannot be repaired by simply replacing a confidence threshold with an information-gain or evidence-strength threshold. The next research question is structural: which representation of recent evidence, temporal context, or model discrepancy can distinguish reliable from unreliable inference before any new reject policy is designed?

\printbibliography

\appendix
\section{Frozen analysis contract}
\label{app:contract}
The text below is the contract as committed in \texttt{34f9797} on 13 September 2026, before the analysis was first executed. It is reproduced verbatim apart from the removal of its cross-reference label, so its tense and references to pending results reflect the manuscript at that point.

\subsection*{Pre-result hypotheses and interpretation contract}

The following diagnostic hypotheses are fixed before the new panel result is inserted into this manuscript:
\begin{enumerate}[label=H\arabic*.,leftmargin=*]
    \item The median household information-gain separation exceeds the no-separation value, $\operatorname{median}_h(A_{I,h})>0.5$.
    \item Information gain separates correctness more strongly than posterior confidence, $\operatorname{median}_h(A_{I,h})>\operatorname{median}_h(A_{C,h})$.
    \item Interval-level evidence strength carries positive correctness information, $\operatorname{median}_h(A_{E,h})>0.5$.
\end{enumerate}

These are descriptive directional hypotheses, not null-hypothesis significance tests. No equivalence margin, significance threshold, or post-hoc cut-off will be introduced after the panel is run. The natural reference point $A=0.5$ and the pre-stated metric ordering are the only decision anchors.

The interpretation is also fixed in advance:
\begin{itemize}
    \item If H1 and H2 hold, posterior concentration and update-level information are empirically distinct, and information gain becomes a candidate quantity for a later selective-prediction study.
    \item If H1 fails, replacing confidence with information gain is not supported by this development panel, regardless of whether information gain is numerically larger than confidence.
    \item If confidence is inverted while information gain is not, the failure is consistent with accumulated posterior structure carrying confidence beyond what the current update justifies.
    \item If both confidence and information gain are weak or inverted, the failure cannot be repaired by exchanging one scalar threshold statistic for the other; the observation/inference formulation itself requires further study.
    \item If H3 holds while H1 fails, raw evidential magnitude contains correctness information that is being transformed poorly by the Bayesian update geometry.
\end{itemize}

The paper will not select a preferred abstention threshold, optimise coverage, or define a new deployment rule from these 22 development homes. Any such policy belongs to a subsequent study with its own validation design.

\subsection*{Pre-specified outputs}

The machine-readable JSON generated by \texttt{scripts/analyse\_uncertainty\_panel.py} is the source of truth for all new empirical values in this paper. The analysis has two pre-specified paper-facing outputs:
\begin{enumerate}[label=(\roman*)]
    \item \cref{tab:separation}, containing the median within-home correctness AUC and contributing-home count for all five diagnostics;
    \item a household-distribution figure showing one point per home for each diagnostic, with a horizontal reference at $A=0.5$ and no fitted trend or threshold line.
\end{enumerate}

The figure is descriptive: it is intended to show heterogeneity and direction across homes, not to convert the development panel into a confirmatory test. If the JSON artifact does not exist, both the table values and the figure remain pending rather than being reconstructed manually from console output.

\end{document}
